%!TEX root = ../Thesis.tex
\chapter{Introduction}
Assessing the impact of climate change and developing early-warning strategies is essential for a sustainable future on planet Earth; every year now marks a new record for the hottest ever recorded, the Earth has more excess energy than ever before, and the Arctic ice caps are melting at a rapid rate \cite{un-climate}. To create such an assessment we must gain a better understanding of the climate, not only at a global scale, but on a local, fine-resolution scale \cite{hclim-fennoscandinavia} \cite{precip-temp-downscaling-cnn} \cite{precip-downscaling} \cite{icyalert-article} \cite{review-ml-downscaling}. 
With a such understanding, we can make more precise models of how the climate will develop over a long time-frame and how that development will change the terms for life on earth. The models will help estimate climate tipping points; like the eradication of the Amazon forest, or the first ice-free winter in the arctic region, or extreme weather events; like heatwaves and cloudbursts. 

The current \gls{sota}-method for creating global scale climate models is \gls{gcm}. They capture the oscillating weather systems of the planet and reveal the general direction of Earth systems. But too zoomed out, they often overlook small scale variability; extreme hydrological events are missed, temperatures are poorly mapped, and local wind phenomena are unresolved \cite{precip-downscaling} \cite{temp-downscaling-rnn}. \Glspl{ddm}, on the other hand, model the climate on a regional or local scale by iteratively \gls{downscaling} distance between the data points of the \glspl{gcm}. They use locally-based physics engines to add precision and clarity. These capture the local variations much better, and are considered the \gls{sota} of climate \gls{downscaling} methods \cite{hclim-fennoscandinavia}. However, these models are limited by their intensive computational costs and are therefore often not viable for large scale inference \cite{precip-temp-downscaling-cnn}. 

Instead, recent approaches have implemented \gls{sdm} which utilize a machine learning approach to develop output similar to that of the \glspl{ddm} \cite{precip-temp-downscaling-cnn} \cite{tomasi} \cite{marie} \cite{dynamical-downscaling} \cite{review-ml-downscaling}. \Glspl{sdm} come with several advantages, the most important of which is a significant increase in inference time and proportional decrease in cost. However, the field of \glspl{sdm} lacks maturity, as much of the current literature relies on non-\gls{sota} computer vision methods, such as ResNets, VAEs and GANs \cite{review-ml-downscaling} \cite{review-ml-downscaling-resnet}. Only recent approaches have researched the field of diffusion-based \gls{downscaling} \cite{marie}, but even then explain the lack of further research in the field \cite{tomasi}. As such, there is a vacuum in the field, which must be filled with recent \gls{sota} diffusion-based image generation techniques for climate \gls{downscaling}.

In this thesis, I propose a diffusion-based model for spatial temperature and precipitation \gls{downscaling} of the the EC Earth \gls{gcm} over the European domain. To improve the model design over previous attempts I will implement recent image generation techniques, such as residual target learning (following the approach of \cite{tomasi}), experimenting with attention blocks (like \cite{marie}), and evaluating the use of ensemble samples for estimating uncertainties. Finally I will evaluate the model output by analysing extreme weather events, such as heat-waves, to make sure the model aptly captures the climate variability. I will target three time periods; 1985-2015, 2020-2050, 2070-2100 and compare the development in frequency and strength of such events with the \gls{ddm}.

The above citations summarize to the following hypothesis:
The diffusion-based \gls{downscaling} approach will have comparable performance to that of the \gls{rcm} over the European Domain (similar to what is seen in \cite{tomasi}). Furthermore, the extreme weather events can be modelled with little loss and the diffusion model will accurately depict the temporal development of such events. \\
The diffusion-based approach will have a higher computational cost than simple \glspl{sdm} (such as ResNets and VAEs), and the modifications made beyond the work of \cite{tomasi} will increase the computational load. However, the diffusion approach will have a lower computational cost than a \gls{ddm} (such as the \gls{hclim} \gls{rcm}). \\
Finally, the diffusion-based approach will have a comparable, but lower, performance when transferred to adjacent \glspl{realisation}. Especially, when the correlation between the \glspl{realisation} of the \glspl{gcm} are lower, the model will yield a reduced accuracy. 

\section*{Research Questions}
To test the above hypothesis, below research questions have been formulated:
\begin{enumerate}
    \item How accurately can a diffusion model downscale temperature and precipitation over the European domain, relative to comparative methods?
    \item What are the modifications to the conditional diffusion model approach (as presented by Ho, et. al. in \cite{ddpm}), which can improve model performance for the climate \gls{downscaling} tasks, and and what improvements and trade-offs do they introduce? 
    \item How well does a diffusion model reproduce extreme weather events; such as heatwaves and cloudbursts, relative to comparative methods?
    \item To what extend can a diffusion model, trained on one \gls{gcm}-\gls{rcm} realization generalise to other \glspl{realisation} of the same climate model pair?
\end{enumerate}